\section{Experiments}
\label{sec:experiments}
\subsection{Experimental Setup}
\label{sec:exp_setup}
\paragraph{Models.}
Main experiments are conducted on two 14B backbones: DeepSeek-R1-Distill-Qwen-14B and Qwen3-14B.
To reduce computational cost, analysis experiments are conducted on Qwen3-4B-Instruct-2507. 
All models are trained using the same pool of 300 verifiable environments described in Section \ref{sec:atomic}.

\paragraph{Evaluation}
Models are evaluated on six benchmarks: LiveCodeBench~\citep{jain2025livecodebench}, AIME 2024/2025~\citep{AIME2024I,AIME2025I}, Enigmata~\citep{chen2026enigmata}, IFEval~\citep{zhou2023instruction}, and LongBench-v2~\citep{bai-etal-2025-longbench}.
These benchmarks encompass code generation, mathematical reasoning, logic reasoning, instruction following, and long-context understanding.
None of these benchmarks are used in RACES. Consequently, improvements on these challenging benchmarks indicate genuine reasoning generalization.
To mitigate test randomness, each benchmark is evaluated multiple times (32 times for AIME 2024/2025 and 4 times for other benchmarks), and the average score is reported.

\paragraph{Data Construction.}
Two configurations sharing the same environment pool are compared: an individual-environment baseline that samples instances directly from the pool, and RACES, which samples a composition operator, a composition size, and domain-compatible environments to form a composite. Composition sizes are uniformly drawn from $[2,12]$ for \textsc{SEQUENTIAL} and \textsc{PARALLEL}, and $[2,6]$ for \textsc{SORT} and \textsc{SELECT}. In the analysis experiments, these ranges are reduced to $[2,6]$ for \textsc{SEQUENTIAL} and \textsc{PARALLEL}, and $[2,3]$ for \textsc{SORT}. An anonymized subset is available at \url{https://anonymous.4open.science/r/Submission_of_NIPS2026_34776-B7FD}.

\paragraph{Training Details.}
All RL experiments are implemented with VERL on a 32 × NVIDIA A100 80GB cluster and vLLM for rollout generation. Optimization is performed with GRPO~\citep{shao2024deepseekmathpushinglimitsmathematical} using a clip ratio of 0.28 and no KL regularization to a reference policy. For the main experiments, each configuration is trained for 300 steps using 12,800 training instances, a batch size of 128, a learning rate of 2e-6, 8 rollouts per problem, and a maximum sequence length of 32K tokens. For the analysis experiments, Qwen3-4B-Instruct-2507 is trained for 200 steps using 6,400 instances, a batch size of 64, and a maximum sequence length of 16K tokens. \textsc{Select} is excluded from the analysis as it yields near-zero rewards on this smaller backbone.

\subsection{Main Results}
\label{sec:main_results}
\begin{table}[!t]
\centering
\caption{\textbf{Main results.} LCBench and LBench-V2 denote LiveCodeBench and LongBench-V2; AIME is the average score on AIME 2024 and AIME 2025. Both RL settings use the same number of training instances and steps.}
\resizebox{0.8\linewidth}{!}{
\begin{tabular}{lcccccc}
\toprule
Model & LCBench & Enigmata & LBench-V2 & IFEval & AIME & Avg. \\
\midrule
\multicolumn{7}{l}
{\textcolor{gray}{\textit{DeepSeek-R1-Distill-Qwen-14B as base model}}}\\
\noalign{\vskip 0.11cm}
Base & 47.2 & 32.3 & 32.5 & 70.6 & 58.5 & 48.2 \\
$RL_{individual}$ & 46.9 & 34.2 & 33.7 & 69.3 & 59.8 & 48.8 \\
$RL_{RACES}$ & \textbf{48.8} & \textbf{35.4} & \textbf{36.0} & \textbf{74.6} & \textbf{61.7} & \textbf{51.3} \\
\noalign{\vskip 0.11cm}
\hdashline
\noalign{\vskip 0.11cm}
\multicolumn{7}{l}
{\textcolor{gray}{\textit{Qwen3-14B as base model}}}\\
\noalign{\vskip 0.11cm}
Base & 55.0 & 47.4 & 32.5 & 84.5 & 74.8 & 58.8 \\
$RL_{individual}$ & 56.3 & 48.2 & 34.1 & 85.7 & 76.0 & 60.1 \\
$RL_{RACES}$ & \textbf{57.0} & \textbf{49.2} & \textbf{35.5} & \textbf{86.7} & \textbf{77.0} & \textbf{61.1} \\
\bottomrule
\end{tabular}
}
\label{tab:main_results}
\end{table}

Table~\ref{tab:main_results} presents a comparison between RL on RACES ($RL_{RACES}$) and RL on individual environments ($RL_{individual}$).
$RL_{RACES}$ and $RL_{individual}$ are initialized from the same environment pool and utilize identical RL data scales.
Across both model families, RACES consistently demonstrates superior performance.
For DeepSeek-R1-Distill-Qwen-14B, $RL_{individual}$ increases the average score marginally, from 48.2 to 48.8.
In contrast, RACES improves the average score to 51.3, outperforming $RL_{individual}$ by 2.5 points.
For Qwen3-14B, RACES also surpasses $RL_{individual}$, increasing the average score from 60.1 to 61.1.
Performance gains are consistent across all benchmarks for both backbones.

$RL_{individual}$ exposes the model to isolated environments, whereas $RL_{RACES}$ trains the model on composed environments that require chaining transformations, maintaining intermediate states, and inferring valid operation orders or subsets.
The consistent advantage of RACES indicates that environmental composition offers a more effective training signal than sampling instances from the original environment pool, resulting in stronger transfer on unseen reasoning benchmarks.